\section{Drivers}

Drivers provide methods for interacting with different hardware components on a device.
Typically, hardware components expose registers that control various functionalities.
The role of a driver is to abstract these low-level register interactions and provide a
user-friendly interface.

\subsection{Timers}

The Nordic nRF51 board on the \texttt{micro:bit} features three timers:
\begin{itemize}
    \item \texttt{TIMER0} – A 32-bit timer.
    \item \texttt{TIMER1} and \texttt{TIMER2} – 16-bit timers.
\end{itemize}

Apollo-RTOS configures the two 16-bit timers as 1MHz timers, generating an interrupt
every 1s. These timers serve distinct purposes:
\begin{itemize}
    \item \lstinline[style=myagc]|TIMER1| – Maintains the system clock.
    \item \lstinline[style=myagc]|TIMER2| – Used for profiling.
\end{itemize}

The initialization function for both timers is straightforward, primarily involving
writing values to specific registers in the correct order. This implementation was
inspired by \citet{microbian}.

The interrupt handling for \lstinline[style=myagc]|TIMER1| (the system timer) is more
intricate and will be discussed next.
\begin{lstlisting}[
    style=myagc,
    title={\file{src/drivers/timer.cpp}},
    linebackgroundcolor={\lines{14}{15} \line{20} \lineD{7}},
]
volatile time_t MILLIS = 0;

extern "C" void timer1_handler(void)
{
  // increment MILLIS

  sched::assert_stack(); // check for stack overflow 

  // run the waitlist tasks every (update_interval)us
  if (MILLIS % waitlist::update_interval == 0) {
    intr_guard guard;

    // setup temporary stk to run waitlist tasks
    auto prev_stk = (uint8_t *)get_msp();
    set_msp(/* temporary stack */); 

    waitlist::run();

    // restore original stack
    set_msp(prev_stk);
  }
}
\end{lstlisting}

As mentioned earlier in \autoref{sec:waitlist}, at regular intervals (every
\lstinline[style=myagc]|update_interval|\(\mu\)s), the waitlist is checked for any
pending tasks. Since the timer interrupt can occur at any point in execution, there is
a risk of stack overflow. To mitigate this, waitlist tasks are executed on a separate
stack. This is why a temporary stack is created and set up in lines 14-15.


\subsubsection{Profiling}

This is an experimental feature currently under development. % FIXME TODO
The idea is to use \lstinline[style=myagc]|TIMER2| to track CPU time consumed by each
process and function. At present, it performs reasonably well for profiling processes
but struggles with function-level profiling. 

For processes, users can type \lstinline[style=mybash]|trace| in the shell to display
process statuses along with their CPU usage: \vspace{-1em}
\begin{center}
\begin{minipage}{.7\linewidth}
\begin{lstlisting}[style=mytxt]
curr_proc: trace
  |  0. idle 	 : (idle),   [runnable], 38.938217%
  |  1. shell 	 : (high),   [asleep],    0.18498%
  |  2. display  : (low),    [runnable], 48.945616%
  |  3. accel_bg : (high),   [runnable],  0.18498%
  |  4. test 	 : (low),    [asleep],    1.609323%
  |  5. trace 	 : (medium), [runnable], 14.95450%
\end{lstlisting}
\end{minipage}
\end{center}
As seen in the output, most CPU time is currently spent in the
\lstinline[style=myagc]|idle| and \lstinline[style=myagc]|display| processes—mainly
because I haven't written a more demanding workload. A significant portion is also used
for printing this information over the serial bus, reaffirming the well-known overhead
of I/O operations.

This behavior is implemented in the \lstinline[style=myagc]|timer2_handler|:
\begin{lstlisting}[
    style=myagc,
    title={\file{src/drivers/timer.cpp}},
    linebackgroundcolor={\lines{7}{8}},
]
volatile time_t debug_ticks = 0;

extern "C" void timer2_handler(void)
{
  // increment debug_ticks

  sched::tick(); // increment tick for the current process
  profile::tick(); // increment tick for the currently running function
  // but only do this for functions that have profiling enabled
}
\end{lstlisting}


\subsection{Serial}

The serial driver manages the UART connection. The
\lstinline[style=myagc]|serial::init()| function initializes the UART device and
enables interrupts. The \lstinline[style=myagc]|serial| namespace provides three core
functions:

\begin{itemize}
    \item \lstinline[style=myagc]|void busy_putc(char)|: A
        \lstinline[style=myagc]|putc| function that disables interrupts and busy-waits
        until the character is sent.
    \item \lstinline[style=myagc]|char getc()|: Waits for a character when interrupts
        are disabled.
    \item \lstinline[style=myagc]|void intr_putc(char)|: A buffered
        \lstinline[style=myagc]|putc| implementation that uses interrupts.
\end{itemize}

The \lstinline[style=myagc]|busy_putc| and \lstinline[style=myagc]|getc| functions are
straightforward: they send or receive a character and wait for the device to become
ready. The \lstinline[style=myagc]|intr_putc| function, however, is more interesting
and is partially inspired by the implementation by \citet{microbian}.

\begin{lstlisting}[
    style=myagc,
    title={\file{src/drivers/serial.cpp}},
    linebackgroundcolor={\lines{12}{13} \line{17}},
]
volatile int txidle; /* Whether UART is idle */
circular_buffer<char, NBUF> buf;

void intr_putc(char ch)
{
  // if the buffer is full
  while (buf.size() == NBUF)
    pause(); // enter standby state waiting for an interrupt

  // if the UART is idle, send the character immediately
  if (txidle) {
    UART.TXD = ch; 
    txidle = 0; // mark the UART busy
  } 
  // otherwise enqueue the character
  else 
    buf.enqueue(ch);
}
\end{lstlisting}

With this setup, when functions like \lstinline[style=myagc]|printf| send multiple
characters, usually only the first character is transmitted immediately. The remaining
characters are placed into a circular buffer, and it becomes the UART interrupt
handler’s responsibility to send them asynchronously.

\begin{lstlisting}[
    style=myagc,
    title={\file{src/drivers/serial.cpp}},
    linebackgroundcolor={\line{9} \line{11}},
]
extern "C" void uart_handler(void)
{
  // ... handle the receive path here

  // UART finished transmission
  if (UART.TXDRDY) { 
    UART.TXDRDY = 0;
    if (buf.empty()) // no more characters to send
      txidle = 1; // mark UART idle
    else
      UART.TXD = buf.dequeue(); // otherwise send the next character
  }

  // ... clear and re-enable interrupts
}
\end{lstlisting}


\subsubsection{Listening on the UART}

In Apollo-RTOS, all incoming UART characters are, by default, sent to the
\lstinline[style=myagc]|shell| process. This creates the illusion that the entire
operating system is simply the shell, allowing user interaction through commands.

However, some processes may need to momentarily take control of input, similar to
starting a Python shell within a Bash shell. To enable this, the
\lstinline[style=myagc]|serial| driver maintains a stack of ``listeners" that process
incoming characters. The UART handler sends each character from the top of the listener
stack downward. If a listener ``consumes" the character, it is not passed to any other
listeners.

To explain what I mean:
\begin{lstlisting}[
    style=myagc,
    title={\file{include/drivers/serial.h}},
]
using listener_t = bool (*)(char);

void register_listener(listener_t);
void unregister_listener();
\end{lstlisting}
\begin{lstlisting}[
    style=myagc,
    title={\file{src/drivers/serial.cpp}},
    % linebackgroundcolor={\line{9} \line{11}},
]
class {
  listener_t stack[16] = {nullptr};
  size_t idx = 0;

public:
  void push(listener_t listener) { stack[idx++] = listener; }
  void pop() { idx--; }

  void send(char c)
  {
    auto i = idx;
    while (i > 0) {
      if (stack[--i](c)) 
        break; // if consumed, don't send c to the rest of the listenrs
    }
  }
} listeners;

void register_listener(listener_t listener) { listeners.push(listener); }
void unregister_listener() { listeners.pop(); }
\end{lstlisting}

Finally in the UART handler, all incoming characters get sent to
\lstinline[style=myagc]|listeners|

\begin{lstlisting}[
    style=myagc,
    title={\file{src/drivers/serial.cpp}},
]
extern "C" void uart_handler(void)
{
  if (UART.RXDRDY) {
    listeners.send(UART.RXD); // send the character to the listeners
    UART.RXDRDY = 0;
  }

  // ... the rest
}
\end{lstlisting}

Of course, this registering and unregistering pattern can be put into an RAII class:
\begin{lstlisting}[
    style=myagc,
    title={\file{include/drivers/serial.h}},
]
struct listener_guard {
  listener_guard(listener_t listener) { register_listener(listener); }
  ~listener_guard() { unregister_listener(); }
};
\end{lstlisting}


\subsection{I2C}

The I2C master on the Nordic nRF51 enables communication with devices connected to the
\texttt{micro:bit} via the I2C bus. It operates at frequencies of 100KHz or 400KHz. For
this project, we use the 100KHz configuration provided by \citet{microbian}.

Apollo-RTOS includes two I2C driver implementations: synchronous and asynchronous. The
synchronous implementation is straightforward, with all read and write operations
busy-waiting until completion. This implementation is used before the scheduler is
initialized.

Once the scheduler is active, the asynchronous implementation takes over. In this mode,
the I2C backend uses a mutex to synchronize requests from different processes. From the
process's perspective, each I2C operation appears atomic. However, due to the nature of
I2C communication, sending multiple bytes triggers multiple interrupts. To maintain the
illusion of atomicity, the process is put to sleep while communication is in progress.

Below is a brief explanation of how the I2C master communicates with slave devices:
\begin{itemize}[left=0pt]
    \item First, the device address is written to the
        \lstinline[style=myagc]|I2C0.ADDRESS| register. This address is unique to each
        device and ranges between \texttt{0x03} and \texttt{0x77}.
    \item Next, the bytes constituting the "command" are written. After each byte, an
        acknowledgment interrupt is raised, signaling the system to send the next byte.
    \item After all command bytes are sent, the process differs for read and write
        operations:
        \begin{itemize}[left=0pt]
            \item For write operations, bytes are sent sequentially, with
                acknowledgment interrupts confirming each transmission.
            \item For read operations:
                \begin{itemize}[left=0pt]
                    \item The \lstinline[style=myagc]|STARTRX| register is set to
                        \texttt{1}.
                        \begin{itemize}[left=0pt]
                            \item If this is the last byte to be read, the
                                \lstinline[style=myagc]|STOP| bit in the
                                \lstinline[style=myagc]|SHORTS| register is set.
                            \item Otherwise, the \lstinline[style=myagc]|SUSPEND| bit
                                is set.
                        \end{itemize}
                    \item The device writes the requested data to the
                        \lstinline[style=myagc]|RXD| register and raises an interrupt.
                    \item The interrupt is detected, and the data is read from the
                        \lstinline[style=myagc]|RXD| register.
                    \item If more bytes are to be read, the
                        \lstinline[style=myagc]|RESUME| register is set to \texttt{1}.
                \end{itemize}
        \end{itemize}
\end{itemize}

To manage this process, we implement a state machine with four states:
\begin{itemize}
    \item \lstinline[style=myagc]|W_CMD|: Writing command bytes.
    \item \lstinline[style=myagc]|TX_DATA|: Transmitting data.
    \item \lstinline[style=myagc]|RX_DATA|: Receiving data.
    \item \lstinline[style=myagc]|NACK|: Finalizing the operation and performing
        cleanup.
\end{itemize}

The communication begins in the \lstinline[style=myagc]|xfer| method:
\begin{lstlisting}[
    style=myagc,
    title={\file{src/drivers/i2c\_async.h}},
    linebackgroundcolor={}
]
template <bool is_read>
int xfer(uint8_t addr,
         uint8_t *cmd, size_t cmd_sz, 
         uint8_t *buf, size_t buf_sz)
{
  // acquire lock before continuing
  mtx.lock(); 

  // ... copy the buffer addresses and buffer sizes

  // initial state is W_CMD
  state.stage = stage_t::W_CMD;
  I2C0.ADDRESS = addr;

  // start transmission and send the first byte
  I2C0.STARTTX = 1;
  I2C0.TXD = *cmd;

  // put the current process to sleep until woken up after completion
  sched::wait(intr_chan);
   
  // ... report any error that might've occured
}
\end{lstlisting}

This method just starts the transmission, and puts the process to sleep. The real magic
(or nightmare if implemented incorrectly) happens in the interrupt handler:

\begin{lstlisting}[
    style=myagc,
    title={\file{src/drivers/i2c\_async.h}},
    linebackgroundcolor={\lines{5}{6} \lines{8}{30} \lines{32}{44} \lines{46}{60}
    \lines{62}{67}}
]
void handler(void)
{
  intr_guard guard{I2C0_IRQ};

  if (I2C0.ERROR)
    // ... if there was an error, cancel the communication

  else if (state.stage == stage_t::W_CMD) {
    // ... receive acklowledgement                                      
                                                         +------------+ 
    // still writing the command                         |  W_CMD     | 
    if (state.cmd_idx < state.cmd_sz)                    +------------+ 
      // ... send next byte                                             

    // no data to send or receive, so stop 
    else if (state.data_sz == 0)
      // ... stop

    // if we're writing
    else if (state.mode == mode_t::WRITE) {
      state.stage = stage_t::TX_DATA;
      // ... send the first data byte
    }

    // we're reading
    else {
      state.stage = stage_t::RX_DATA;
      // ... set the SHORTS register and STARTRX = 1
    }
  }

  else if (state.stage == stage_t::TX_DATA) {
    // ... receive acklowledgement
                                                         +------------+              
    // finished writing                                  |  TX_DATA   |               
    if (state.data_idx == state.data_sz)                 +------------+               
      // ... stop                                                       

    // more bytes to write
    else {
      state.stage = stage_t::TX_DATA;
      // ... send the next byte
    }
  }

  else if (state.stage == stage_t::RX_DATA) {
    // ... receive acklowledgement
                                                         +------------+              
    // ... copy byte from I2C0.RXD                       |  RX_DATA   |               
                                                         +------------+               
    // finished reading, move to NACK
    if (state.data_idx == state.data_sz)
      state.stage = stage_t::NACK;

    // continue reading
    else {
      state.stage = stage_t::RX_DATA;
      // ... set the SHORTS register and RESUME = 1
    }
  }

  else if (state.stage == stage_t::NACK) {               
    // ... receive acklowledgement                      +------------+
                                                        |  NACK      | 
    mtx.unlock(); // release the lock                   +------------+
    sched::notify(intr_chan); // then wake up the waiting process
  }
}
\end{lstlisting}

With this implementation, multiple processes can communicate over the I2C bus without
conflicts. All requests are serialized using the \lstinline[style=myagc]|mutex|,
ensuring orderly access to the bus. Additionally, while one process is communicating
over the I2C bus, it does not monopolize the CPU. Instead, it allows other processes to
continue executing, improving overall system efficiency.




\subsection{Display}

\subsection{Display}

The \(5 \times 5\) LED matrix is arranged into three columns in a counter-intuitive
format. To control these LEDs, we utilized the convenience macros provided by
\citet{microbian}. The LEDs are managed through GPIO registers and require constant
updates to display smooth animations. To achieve this, the
\lstinline[style=myagc]|display| driver is implemented as a low-priority process, with
its priority set just above the \lstinline[style=myagc]|idle_proc|.

The implementation is straightforward:
\begin{lstlisting}[
    style=myagc,
    title={\file{src/drivers/display.cpp}},
    linebackgroundcolor={\line{12}}
]
image_t image;

/* definition of display */
{
  GPIO.DIR = LED_MASK;
  reset(); // clear the image out at the beginning

  while (true) {
    int n = 0;
    // display the three rows of the image, and wait for a little
    while (n < 3) {
      GPIO.OUT = image[n++];
      delay_loop(5000);
    }
    sched::sleep(10);
  }
}
\end{lstlisting}

A method named \lstinline[style=myagc]|set_image| is provided to allow processes to
update the global variable \lstinline[style=myagc]|image| with the desired display
image. The display driver continuously renders the image stored in this variable.

As discussed earlier in \autoref{sec:rec_init}, this process is always started during
system initialization.


